\documentclass[11pt]{article}
\usepackage[margin=1in]{geometry}
\usepackage{booktabs}
\usepackage{amsmath}
\usepackage{hyperref}
\title{Analytic Partial-Residual Targets: A Controlled Mechanism Study and its Structural Boundary}
\author{Anonymous authors}
\date{Draft generated from recorded experiment artifacts}

\begin{document}
\maketitle

\begin{abstract}
We study an exact partial-residual target update for grouped additive linear regression.  The update is an instrumented form of block-coordinate minimization/backfitting, rather than a learned target generator.  In the registered S1 synthetic additive matrix, it met the two registered endpoints against development-selected Adam on five final seeds, but failed the quality endpoint on a separate longer-horizon family.  We then registered two structural-misspecification successors with interleaved supplied groups and unmodeled nonlinear/interacting targets.  S2.1 adds a common affine intercept after S2 exposed that the initial nonlinear family had a mean-offset confound.  S2.1 failed both registered endpoints on its final and long-horizon families.  These internally executed synthetic findings establish a narrow boundary, not general learning superiority.
\end{abstract}

\section{Question and mechanism}
For a grouped additive model $f(x)=\sum_g f_g(x_g)$, the proposed update constructs $t_g=y-\sum_{h\ne g} f_h(x_h)$ and solves a ridge subproblem for group $g$.  The method uses support labels and current predictions only.  It has no teacher, learned inverse, component labels, or trained target controller.  Consequently, its most accurate interpretation is partial-residual backfitting/block-coordinate minimization \cite{wright2015,buja1989}, not a general target-propagation result \cite{lee2015,jaderberg2016}.

\section{Protocols and evidence boundary}
S1 freezes selection on meta-training and development task families, with final and long-horizon labels excluded from selection.  It compares tuned full-batch Adam, sequential analytic residual updates, a no-derivation block-gradient ablation, and fixed normalized LMS under deterministic operation-proxy accounting.  S2 and S2.1 use distinct versioned protocols and seeds.  S2.1 keeps the same information boundary but adds a common learned intercept to all conditions.

The tables and figures are generated from saved outputs and independently read back; no numeric result is entered by hand.  The exact commands, configuration hashes, and manifests are retained alongside each run.

\section{Results}
\begin{table}[t]
\centering
\caption{Recorded equal-resource evaluation NMSE.  ``Fail'' means the predeclared composite endpoint was not supported; it does not mean a runtime error.}
\begin{tabular}{llrrl}
\toprule
Study & Split & Adam & Analytic residual & Registered decision \\
\midrule
S1 & final & 0.024140 & 0.003475 & supported on admitted matrix \\
S1 & long horizon & 0.004701 & 0.004754 & fail \\
S2.1 & final & 0.724612 & 0.711624 & fail \\
S2.1 & long horizon & 0.822250 & 0.822343 & fail \\
\bottomrule
\end{tabular}
\label{tab:primary-results}
\end{table}

\begin{table}[t]
\centering
\caption{S2.1 descriptive paired result based on the five frozen task pairs per split.  Intervals are deterministic percentile-bootstrap intervals for the mean paired NMSE difference (analytic residual minus Adam).}
\begin{tabular}{lrrr}
\toprule
Split & analytic wins & mean difference & 95\% interval \\
\midrule
final & 3/5 & -0.0129873 & [-0.0281100, 0.0021354] \\
long horizon & 1/5 & 0.0000932 & [-0.0003178, 0.0004638] \\
\bottomrule
\end{tabular}
\label{tab:s21-paired}
\end{table}


In S1, the analytic method passed the five-seed final matrix but the longer-horizon matrix failed the registered quality criterion.  Thus even the original additive result does not support a persistent long-budget advantage.  S2's structural matrix was completed and verified, but its nonlinear target had a nonzero mean while no model fitted an intercept; it is retained as a negative run rather than used as a clean structural attribution.  S2.1 corrects that common-capacity omission prospectively.  It fails the registered quality and efficiency criteria in both final and long-horizon splits.  Its paired bootstrap intervals are descriptive only: the protocol endpoints, not post-hoc intervals, control the claim decision.

\section{Limitations and negative evidence}
The work remains synthetic, CPU-only, internally executed, and small-seed.  It has no independent reproduction, no real-data admission track, no neural block implementation, no learned controller, and no wall-clock or memory comparison under a separately frozen protocol.  The simultaneous/Jacobi residual ablation is numerically unstable in S2/S2.1 under the selected settings; this is retained as a result, not suppressed.  The paper therefore makes no publication-ready claim of general learning utility, deep credit assignment, or learned-target superiority.

\section{Reproducibility}
Run \texttt{sh scripts/run\_s21\_reproduction.sh OUTPUT\_DIRECTORY} from the repository root, then inspect the generated manifest and evaluation-verification receipts.  The existing S1 and S2 artifacts are preserved under \texttt{runs/}.  The external reproduction packet remains unfinished.

\bibliographystyle{plain}
\bibliography{references}
\end{document}
